%==============================================================================
\section{Hồi quy tuyến tính đa biến (Multiple Linear Regression)}
\label{sec:multiple-linear-regression}
%==============================================================================

\begin{dinhnghia}[title={Hồi quy tuyến tính đa biến (MLR)}]
Thuật toán học có giám sát mô hình quan hệ giữa \textbf{một} biến phụ thuộc liên tục và \textbf{nhiều} biến độc lập.
Biến phân loại (ví dụ \texttt{State}) cần mã hóa số trước khi đưa vào mô hình tuyến tính.
\end{dinhnghia}

\begin{ynghia}[title={So với hồi quy đơn}]
MLR là mở rộng của hồi quy đơn: dự đoán bằng hai hoặc nhiều feature.
\end{ynghia}

\begin{dinhnghia}[title={Phương trình}]
Với $n$ quan sát, $p$ feature:
\[
h(x_i) = w_0 + w_1 x_{i1} + w_2 x_{i2} + \cdots + w_p x_{ip}
\]
Có phần dư:
\[
y_i = h(x_i) + e_i, \quad e_i = y_i - h(x_i)
\]
\end{dinhnghia}

\subsection{Giả định}

\begin{luuy}[title={Giả định MLR}]
\begin{enumerate}
  \item Tuyến tính giữa target và predictor.
  \item Các quan sát độc lập.
  \item Phương sai phần dư đồng nhất (homoscedasticity).
  \item Phần dư gần phân phối chuẩn.
  \item Không đa cộng tuyến mạnh giữa feature.
  \item Không tự tương quan phần dư.
  \item Giá trị biến độc lập cố định giữa các mẫu lặp.
\end{enumerate}
\end{luuy}

\subsection{Triển khai Python}

\begin{vidu}[title={Dataset data.csv}]
50 dòng: \texttt{R\&D Spend}, \texttt{Administration}, \texttt{Marketing Spend}, \texttt{State}, \texttt{Profit}.
Ví dụ dưới bỏ \texttt{State} như bài gốc (đơn giản hóa).
\end{vidu}

\begin{vidu}[title={Bước 1 --- Chuẩn bị}]
\begin{lstlisting}
import numpy as np
import pandas as pd
import matplotlib.pyplot as plt

dataset = pd.read_csv('data/data.csv')
X = dataset[['R&D Spend', 'Administration', 'Marketing Spend']]
y = dataset['Profit']
print(X.head())
print(y.head())
\end{lstlisting}
\end{vidu}

\begin{vidu}[title={Chia train/test (20\% test)}]
\begin{lstlisting}
from sklearn.model_selection import train_test_split
X_train, X_test, y_train, y_test = train_test_split(
    X, y, test_size=0.2, random_state=0)
\end{lstlisting}
\end{vidu}

\begin{vidu}[title={Bước 2 --- Huấn luyện}]
\begin{lstlisting}
from sklearn.linear_model import LinearRegression
regressor = LinearRegression()
regressor.fit(X_train, y_train)
\end{lstlisting}
\end{vidu}

\begin{vidu}[title={Bước 3 --- Kiểm tra}]
\begin{lstlisting}
y_pred = regressor.predict(X_test)
df = pd.DataFrame({'Real Values': y_test.values,
                   'Predicted Values': y_pred})
print(df)
\end{lstlisting}
\end{vidu}

\begin{vidu}[title={Bước 4 --- Metric}]
\begin{lstlisting}
from sklearn.metrics import mean_squared_error, mean_absolute_error, r2_score
import numpy as np

mse = mean_squared_error(y_test, y_pred)
rmse = np.sqrt(mse)
mae = mean_absolute_error(y_test, y_pred)
r2 = r2_score(y_test, y_pred)
print("MSE:", mse)
print("RMSE:", rmse)
print("MAE:", mae)
print("R2:", r2)
\end{lstlisting}
\end{vidu}

\begin{ynghia}[title={Đọc $R^2 \approx 0{,}96$ (bài gốc)}]
Khoảng 96\% biến thiên của \texttt{Profit} được giải thích bởi ba feature đầu vào (trên tập kiểm tra của ví dụ gốc).
\end{ynghia}

\begin{vidu}[title={Bước 5 --- Dự đoán điểm mới}]
\begin{lstlisting}
new_data = [[166343.2, 136787.8, 461724.1]]
profit = regressor.predict(new_data)
print(profit)
\end{lstlisting}
\end{vidu}

\begin{vidu}[title={Hệ số và chặn}]
\begin{lstlisting}
print("coefficients:", regressor.coef_)
print("intercept:", regressor.intercept_)
\end{lstlisting}
Kiểm tra tay: $\mathrm{Profit} \approx w_0 + w_1 X_1 + w_2 X_2 + w_3 X_3$.
\end{vidu}

\subsection{Ứng dụng và thách thức}

\begin{phuongphap}[title={Ứng dụng}]
Tài chính, marketing, bất động sản, y tế, kinh tế, khoa học xã hội.
\end{phuongphap}

\begin{luuy}[title={Thách thức}]
Đa cộng tuyến, overfitting/underfitting, phi tuyến, outlier, dữ liệu thiếu.
\end{luuy}

\subsection{So sánh đơn và đa biến}

\begin{center}
\small
\begin{tabular}{|l|p{4.2cm}|p{4.2cm}|}
\hline
\textbf{Tiêu chí} & \textbf{Đơn biến} & \textbf{Đa biến} \\
\hline
Số biến độc lập & 1 & $\geq 2$ \\
Phương trình & $y = w_0 + w_1 x$ & $y = w_0 + \sum w_j x_j$ \\
Độ phức tạp & Thấp hơn & Cao hơn \\
Giải thích hệ số & Dễ hơn & Khó hơn \\
\hline
\end{tabular}
\end{center}
